\begin{Solbook}{1}
 \par \noindent  a) $x=3,x=-2$ \hspace {2cm} b) $\left (\dfrac {-7}{2};-3\right )$  \hfill \qedEX 
\end{Solbook}
\begin{Solbook}{2}
 \par \noindent  a) 6 \hspace {2cm} b) $8+3\sqrt {3}$  \hfill \qedEX 
\end{Solbook}
\begin{Solbook}{3}
 \par \noindent  a) $V\approx 105$ (volt) \hspace {2cm} b) $P=\dfrac {275}{2}>100$  \hfill \qedEX 
\end{Solbook}
\begin{Solbook}{4}
 \par \noindent  Nhà trường thuê 7 xe 45 chỗ và 3 xe 16 chỗ.  \hfill \qedEX 
\end{Solbook}
\begin{Solbook}{5}
 \par \noindent  a) $C\approx 2024$ (mm) \hspace {2cm} b) $10\cdot 3\cdot 650\pi =19500\pi \,(\text {mm})\approx 61\,(\text {m})$  \hfill \qedEX 
\end{Solbook}
\begin{Solbook}{6}
 \par \noindent  \immini { \begin {enumerate} \item $A,P,O,M$ cùng nằm trên đường tròn đường kính $OP$ \par Ta có $OP\perp AM, BM\perp AM$ suy ra $OP\parallel BM$. \item $\triangle AOP = \triangle OBN \Rightarrow OP = BN$ \par Ta lại có $BN \parallel OP$ nên $OBNP$ là hình bình hành. \par \item Chứng minh $I$ là trực tâm $\triangle POJ$ suy ra $IJ \perp OP$ (1) \par Chứng minh $K$ là trung điểm $OP$. \par Lại có: $\widehat {APO} = \widehat {OPI} = \widehat {IOP}$ \par Suy ra $\triangle IPO$ cân tại $I$ suy ra $IK \perp OP$ (2) \par Từ (1) và (2) suy ra $I, J, K$ thẳng hàng. \end {enumerate} } { \begin {tikzpicture}[scale=1, font=\footnotesize , line join =round, line cap=round, >=stealth] \def \r {1.8} \path (0,0) coordinate (O) (180:\r ) coordinate (A) (0:\r ) coordinate (B) (-\r ,2.5*\r ) coordinate (x)node[left]{$x$} ($(A)!4/5!(x)$) coordinate (P) ($(O)!(A)!(P)$) coordinate (H) ($(A)!2!(H)$) coordinate (M) ($(O)!1cm!-90:(A)$) coordinate (n') (intersection of B--M and O--n') coordinate (N) (intersection of A--N and O--P) coordinate (K) (intersection of O--N and M--P) coordinate (I) (intersection of O--M and N--P) coordinate (J) ; \draw (O) circle (\r ) (x)--(A)--(B)--(N)--(A)--(M)--(P)--(O)--(N) (P)--(J)--(O) ; \draw [dashed] (K)--(J); \foreach \x /\g in{O/-90,A/180,B/0,P/180,K/160,M/0,N/90,I/120,J/40} \draw [fill=white] (\x ) circle (1.5pt) +(\g :.3) node{$\x $}; \end {tikzpicture} }  \hfill \qedEX 
\end{Solbook}
